% =====================================================================
%  06 模型的评价及推广(检验、灵敏度分析、优缺点)
%  结构:6.1 模型检验 → 6.2 灵敏度分析(指向问题五) → 6.3 优缺点
%
%  占位符约定:\textbf{(待回填:……)} 是给编程手留的数值空位,
%  求解脚本跑完后逐个替换,不要删除。
% =====================================================================
\section{模型的评价及推广}

\subsection{模型检验}

\subsubsection{代理模型精度检验}

问题二建立的 GPR 代理模型采用留出法(训练集与测试集按 $8:2$ 划分,
独立测试样本不参与训练)评估拟合与预测精度,三项性能指标的检验结果
见表~\ref{tab:model-check}。由表可见,三个模型的决定系数 $R^2$ 均
接近于 1,均方根误差与平均绝对百分比误差均处于极低量级,表明代理模型
对仿真数据的拟合与预测精度高,可代替 CFD 求解器用于后续的多目标优化
与不确定性传播分析。

% 待回填:以下数字由 code/src/solve_q2.py 运行输出
\begin{table}[h]
\centering
\caption{GPR 代理模型精度检验结果(测试集)}
\label{tab:model-check}
\begin{tabular}{lccc}
\toprule
性能指标 & $R^2$ & $RMSE$ & $MAPE$ \\
\midrule
热阻 $R_{th}$      & \textbf{(待回填)} & \textbf{(待回填)} & \textbf{(待回填)} \\
压降 $\Delta P$    & \textbf{(待回填)} & \textbf{(待回填)} & \textbf{(待回填)} \\
温度非均匀性 $\Delta T$ & \textbf{(待回填)} & \textbf{(待回填)} & \textbf{(待回填)} \\
\bottomrule
\end{tabular}
\end{table}

\subsubsection{机理规律与仿真数据的一致性检验}

问题一推导的参数影响规律与附件 2 样本数据的控制变量分析结果对照:
针肋宽度比对热阻的影响在样本范围内呈"先降后升"的非单调变化,各组
控制变量下极小值均出现在 $\alpha \approx 0.2$ 附近;歧管深高比增大
使流量分配更均匀,温度非均匀性总体下降且影响幅度较小;针肋排数增大
使压降总体升高,且增幅随宽度比增大而加大,均与机理推导一致,表明机理
模型对系统行为的刻画符合仿真规律。定量吻合程度(平均相对偏差等)待
求解结果回填。

\subsubsection{优化结果的可行性检验}

将问题三所得综合最优方案回代 GPR 代理模型验证目标函数值,与优化
过程中预测值的相对偏差小于 \textbf{(待回填:偏差百分比)} $\%$,
验证了优化结果的可行性与自洽性。

\subsection{灵敏度分析}

结构参数与工况参数对性能指标的敏感性已在问题五中系统完成:结构参数
扰动对三项指标的影响由训练好的代理模型严格刻画,采用 Sobol 全局敏感
性指数(一阶主效应、总效应)与局部归一化敏感性系数定量评估,并基于
蒙特卡洛不确定性传播给出变异系数、$95\%$ 置信上限与越界失效概率;
工况参数(流量、入口温度、发热功率)的影响基于标称设计点处的梯度
线性化近似评估,详见 5.5 节。分析表明,系统对 \textbf{(待回填:最敏感
参数与结论)} 的扰动最为敏感,该结论为加工公差控制与运行工况约束提供
了定量依据。

\subsection{模型的优缺点}

\textbf{优点:}
\begin{itemize}
  \item 机理与数据驱动融合:问题一建立 CFD 机理模型,问题二以 GPR
        代理模型逼近高维非线性映射,兼顾物理可解释性与拟合精度;
  \item GPR 除给出点预测外还提供预测方差,为后续优化与不确定性
        传播提供了天然的概率框架;
  \item 采用 NSGA-II 全局搜索与 L-BFGS-B 局部精化结合的混合优化
        策略,并对针肋排数做整数离散修正,兼顾解的全局性与收敛精度;
  \item 决策层完整:理想点法消除量纲差异给出唯一综合方案,Minimax
        后悔值准则降低方案对偏好权重的依赖,变异系数进一步量化了
        方案的稳定性;
  \item 敏感性分析体系完备:局部(LSA)与全局(Sobol)结合,并给出
        越界失效概率等工程可操作指标,直接服务于制造公差与运行
        工况约束设计。
\end{itemize}

\textbf{缺点:}
\begin{itemize}
  \item GPR 代理模型的预测精度依赖训练样本的网格覆盖范围,外推至
        样本空间之外或极端工况时精度难以保证;
  \item 机理模型作了均匀体热源、忽略热辐射等简化假设,与实际芯片
        封装结构的传热过程存在一定偏差;
  \item 蒙特卡洛不确定性传播的收敛性依赖采样规模,参数维度升高时
        计算代价增大;
  \item 问题五中工况参数(流量、入口温度、发热功率)对性能的影响基于
        标称设计点处的梯度线性外推近似,仅在 $\pm 5\%$ 小扰动范围内
        有效,扰动幅度增大时误差随之增大;
  \item Minimax 后悔值准则偏保守,在极端偏好场景下综合性能可能
        略逊于针对性优化方案。
\end{itemize}
